\subsection*{21. Affordability Is Not Willingness to Pay}

Someone may believe a professional program is worth ₹80,000 yet be unable to pay without unsafe borrowing. Another might earn enough to afford it easily but value self-study more. A 3rd may have employer reimbursement but only for accredited offerings. These people occupy different positions in the space of \emph{value}, \emph{liquidity}, and \emph{purchase permission}.

Distinguish 6 distinct dimensions: (a) reservation value, (b) immediately available funds, (c) acceptable financing arrangements, (d) annual or monthly budget, (e) time availability, and (f) the ethical limits of what the founder ought to sell. The viable purchasing region is their intersection, not their average.

For social enterprises and education, this gap creates legitimate business-model opportunities: employer sponsorship, college procurement, scholarships funded from genuine surplus, lower-cost unbundled experiences, or philanthropic support. These options have costs and constraints. Financing does not manufacture ability to repay; poorly designed installment plans can simply move risk onto the learner.

\subsection*{22. The B2B Buyer Is Often a Voting Body}

In an organization, users, champions, executives, finance, procurement, security, and legal can have different reservation values. 1 manager may see ₹10 lakh of workflow benefit, while an IT reviewer sees only vendor risk. The procurement officer may demand 3 quotes even when the preferred solution is clear. The CFO may believe the modeled savings cannot be converted into cash or measurable capacity.

This produces a \emph{collective} willingness-to-pay constraint. It is insufficient that 1 stakeholder loves the product. The vendor needs a credible benefit to the economic buyer, acceptable risk to veto holders, and a feasible budget path. This is why decision packets, pilots, implementation obligations, and transparent measurement can increase achievable contract value without increasing product features.

A useful account map specifies who experiences the pain, who can champion, who owns the budget, who must implement, who can block, and who signs. In complex sales, ``the buyer is willing to pay'' is often a statement about whether this coalition can coordinate: not 1 person's emotional intensity.

\subsection*{23. A Richer Model: WTP as a Moving Frontier}

Imagine a graph with 2 axes: \emph{expected incremental value} and \emph{confidence that the provider will deliver}. Different offers sit in different regions. Increasing technical quality can move the offer right; credible demonstrations can move it upward; reduced switching cost can shift the adoption boundary; new competitors may move the relevant alternative. Buyers with different liquidity constraints face different feasible areas.

This is a useful way to diagnose low conversion. If customers understand and believe the offer but the benefit is small, improve the job selection or product. If expected benefit is large but confidence low, provide proof and bounded trials. If value and confidence are high but payment is impossible, change payer or package. If all 3 are favorable but buyers still postpone, investigate timing, urgency, and internal approval.

There is no universal action called ``increase perceived value.'' Perceived value has causal components. The founder's task is to identify the binding constraint and change the smallest part of the system that can move it.
